\documentclass[10pt,a4paper]{amsart}
\usepackage[margin=19mm]{geometry}
\usepackage{amsmath,amssymb,mathtools}
\usepackage[scr=rsfs,cal=euler]{mathalfa}
\usepackage{xfrac}
\usepackage[unicode,hidelinks,bookmarks=false]{hyperref}
\newcommand{\ie}{i.e.\ }
\newcommand{\cf}{cf.\ }
\newcommand{\resp}{respectively\ }
\newcommand{\Subsection}{\subsection}
\providecommand{\nobreakdash}{\nobreak\hbox{-}}
\setlength{\emergencystretch}{2em}
\multlinegap0pt
\usepackage{amssymb}
\usepackage{mathtools}
\newtheorem{theorem}{Theorem}
\newtheorem{lemma}{Lemma}
\newtheorem{corollary}{Corollary}
\DeclareMathOperator{\Sym}{Sym}
\DeclareMathOperator{\cE}{\mathcal{E}}
\DeclareMathOperator{\cO}{\mathcal{O}}
\title{Braid and Phantom}
\author{Jenia Tevelev}
\date{Source: arXiv:2304.01825v4 (CC BY 4.0)}
\begin{document}
\maketitle

\noindent\emph{Source and licence.} Braid and Phantom. Version arXiv:2304.01825v4 (CC BY 4.0), distributed by arXiv under the Creative Commons Attribution 4.0 International licence, http://creativecommons.org/licenses/by/4.0/, as recorded in the arXiv licence metadata for this exact version and shown on the abstract page https://arxiv.org/abs/2304.01825v4. Full licence notice: \texttt{licenses/arxiv-2304.01825-CC-BY-4.0.txt}.\par\smallskip

\noindent\textit{Editorial scope.} Counted unit: the article's theorem MainTheorem (source0.tex lines 104--128). The article's complete original proof of it is delivered with it (source0.tex lines 3604--3911, one \texttt{proof} environment of 308 lines, which the article places away from the statement). No mathematics is written, repaired or re-derived: the statement and the proof are the article's own lines, in the article's own notation, and no retained line is substituted or re-flowed.  Retained uncounted as the source-local context the proof needs: lines 3165--3183; lines 3259--3266; lines 3323--3327; lines 3514--3523; lines 3560--3582.  Nothing was omitted from any retained span, so the package carries no omission record.  No retained line contains a citation, so the wrapper carries no bibliography and no bibliography entry.  No retained line contains a figure environment, an image-inclusion command or drawing code, so no image is delivered and the package declares no asset dependency.  Editorial formatting only, in addition: an AMS \texttt{amsart} preamble that restates the article's own declarations and macros used by the retained lines, namely the environments \texttt{theorem}, \texttt{lemma}, \texttt{corollary} on separate counters, together with the standard packages those lines need.  The retained environments print in this document as Theorem 1, Lemma 1, Corollary 1 in order of appearance, whereas the article prints the same items with the numbering of its own full document.  The complete native source of the article as distributed in the arXiv e-print for this exact version is bundled unchanged as \texttt{sources/139676/source0.tex}.
\begin{theorem}\label{LoomTheorem}
$D^b(M)$ has a semiorthogonal decomposition whose blocks are arranged into
the following four mega-blocks:
$$
\bigl\langle 
Z^{\lambda+2-g}{\theta^*}^{\lambda-k+1}\cE^{\boxtimes \lambda-2k}
\bigr\rangle_{0\le \lambda\le g-2 \atop 0\le k\le \lfloor{\lambda\over2}\rfloor},\bigl\langle 
Z^{\lambda+3-g}{\theta^*}^{\lambda-k+1}\cE^{\boxtimes \lambda-2k}
\bigr\rangle_{0\le \lambda\le 2(g-2) \atop 0\le k\le  \lfloor{\lambda\over2}\rfloor, \lambda-k\le g-2},$$
$$\bigl\langle 
Z^{\lambda+2-g}{\theta^*}^{\lambda-k}\cE^{\boxtimes \lambda-2k}
\bigr\rangle\!\!\!{}_{0\le \lambda\le 2(g-2) \atop 0\le k\le \lfloor{\lambda\over2}\rfloor, \lambda-k\le g-2}\!\!\!\!,
\bigl\langle 
Z^{\lambda+1-g}{\theta^*}^{\lambda-k-1}\cE^{\boxtimes \lambda-2k}
\bigr\rangle\!\!\!\!{}_{g-1\le \lambda\le 2(g-1) \atop 0\le k\le  \lfloor{\lambda\over2}\rfloor, \lambda-k\le g-1}\!\!\!\!.
$$ 
Within each mega-block, the blocks are arranged in decreasing order of $\lambda$
and those with identical $\lambda$ are further arranged by decreasing $k$.
\end{theorem}

\begin{lemma}\label{pw:direct-images}
Let $p_q:\Sym^qC\times N\to N$ be the projection.  For every $q\geq0$,
$$ R\zeta_*Z^{-q}\simeq
       \bigl(Rp_{q*}\cE^{\boxtimes q}\bigr)^\vee,
 \quad
 R\zeta_*Z^{q+1}\simeq Rp_{q*}\cE^{\boxtimes q}.
$$In particular, $R\zeta_*\cO_M\simeq R\zeta_*Z\simeq\cO_N$.
\end{lemma}

\begin{equation}\label{pw:prefix}
 \mathcal B_{H-2m}(c+m)=
 \left\langle\theta^{c+j}\cE^{\boxtimes(H-2j)}
 \right\rangle_{j=\lfloor H/2\rfloor,\ldots,m}.
\end{equation}

\begin{align}
 \mathbf L_{L\zeta^*\mathcal B_{n-s}(c+s)}
       \langle Z^s\theta^c\cE^{\boxtimes n}\rangle
       &\subset\mathcal A
       &&(n-g+1\le s\le0),\label{pw:left-mutation}\\
 \mathbf R_{L\zeta^*\mathcal B_{n+s}(c)}
       \langle Z^s\theta^c\cE^{\boxtimes n}\rangle
       &\subset\mathcal A'
       &&(0\le s\le g-1-n).\label{pw:right-mutation}
\end{align}

\begin{corollary}\label{pw:available-prefix}
Let $t\in\mathbb Z$ and let $\mathcal B,\mathcal D\subset D^b(N)$ be
admissible. Suppose $Z^tL\zeta^*\mathcal B$ and
$Z^tL\zeta^*\mathcal D$ occur consecutively, in this order, in a
semiorthogonal decomposition of $D^b(M)$.
\begin{enumerate}
\item If an adjacent block $\mathcal C$ follows them and
$R\zeta_*(Z^{-t}\mathcal C)\subset\mathcal B$, then
\[
 \mathbf L_{Z^tL\zeta^*\langle\mathcal B,\mathcal D\rangle}\mathcal C
 =\mathbf L_{Z^tL\zeta^*\mathcal B}\mathcal C
 \subset Z^t\mathcal A.
\]
\item If an adjacent block $\mathcal C$ precedes them and
$R\zeta_*(Z^{1-t}\mathcal C)\subset\mathcal B$, then
\[
 \mathbf R_{Z^tL\zeta^*\langle\mathcal B,\mathcal D\rangle}\mathcal C
 =\mathbf R_{Z^tL\zeta^*\mathcal B}\mathcal C
 \subset Z^t\mathcal A'.
\]
\end{enumerate}
In both cases, the blocks in the family retain their embeddings and order.
\end{corollary}

\begin{theorem}\label{MainTheorem}
$D^b(N)$ has a semiorthogonal decomposition into blocks equivalent to
$D^b(\operatorname{Sym}^kC)$: two blocks for $0\le k\le g-2$ and one
block for $k=g-1$. These blocks are
embedded by the following Fourier--Mukai functors,
$$
\Bigl\langle
\bigl\langle 
{\theta}^{1-g+k}\otimes\cE^{\boxtimes g-2-2k}
\bigr\rangle_{0\le k\le \lfloor{g-2\over2}\rfloor},\quad
\bigl\langle 
{\theta}^{2-g+k}\otimes\cE^{\boxtimes g-3-2k}
\bigr\rangle_{0\le k\le  \lfloor{g-3\over2}\rfloor},\quad{}$$
$${}\quad\bigl\langle 
{\theta}^{2-g+k}\otimes\cE^{\boxtimes g-2-2k}
\bigr\rangle_{0\le k\le \lfloor{g-2\over2}\rfloor},
\quad \bigl\langle 
{\theta}^{2-g+k}\otimes\cE^{\boxtimes g-1-2k}
\bigr\rangle_{0\le k\le  \lfloor{g-1\over2}\rfloor}
\Bigr\rangle,
$$ 
arranged into four mega-blocks.
Within the mega-blocks, the 
blocks are arranged in decreasing order of~$k$.
\end{theorem}

\begin{proof}[Proof of Theorem~\ref{MainTheorem}]
Start with the  decomposition of Theorem~\ref{LoomTheorem}
ordered by decreasing $\lambda$ and then by decreasing $k$.  
Split the megablock II at
$\lambda=g-2$ into IIa and IIb and the megablock III at $\lambda=g-2$ into IIIa and IIIb: the mega-blocks
IIa and IIIa have $\lambda\geq g-2$, while IIb and IIIb have
$\lambda\leq g-3$.  We process the mega-blocks in the order
IV, I, IIa, IIb, IIIb, IIIa.
For $g=2$, the parts IIb and IIIb are empty and their steps below are omitted.

In every leftward procedure take the leftmost block not yet processed;
in every rightward procedure take the rightmost such block.
With this order, each block to which Corollary~\ref{pw:available-prefix}
is applied is adjacent to the retained family. The embeddings and
orders of the retained blocks remain unchanged.


\smallskip
\noindent\textit{Mega-block IV.}
Its blocks are
\begin{equation}\label{sRGSRGDRH}
 \left\langle Z^{\lambda+1-g}\theta^{m+1-\lambda}
       \cE^{\boxtimes(\lambda-2m)}\right\rangle_{{g-1\leq\lambda\leq2(g-1)}\atop
       {0\leq m\leq\lfloor\lambda/2\rfloor,\ \lambda-m\leq g-1}}.
\end{equation}
Retain the  family of rightmost blocks with $\lambda=g-1$:
\begin{equation}\label{dthehwe5jwjw6}
 \left\langle\theta^{2-g+k}\cE^{\boxtimes(g-1-2k)}
 \right\rangle_{k=\lfloor(g-1)/2\rfloor,\ldots,0}
       =L\zeta^*\mathcal B_{g-1}(2-g).
\end{equation}

\begin{lemma}\label{megablockiv}
Let $g\leq\lambda\leq2(g-1)$,
$0\leq m\leq\lfloor\lambda/2\rfloor$, and $\lambda-m\leq g-1$.
Put $k_0=g-1-\lambda+m$ and let $\mathcal B\subset D^b(M)$ be the
initial segment of \eqref{dthehwe5jwjw6} with
$
 k=\lfloor(g-1)/2\rfloor,\ldots,k_0$.
Then
\[
 \mathbf R_{\mathcal B}
 \left\langle Z^{\lambda+1-g}\theta^{m+1-\lambda}
       \cE^{\boxtimes(\lambda-2m)}\right\rangle
       \subset\mathcal A'.
\]
\end{lemma}

\begin{proof}
Set $n=\lambda-2m$, $s=\lambda+1-g$, and $c=m+1-\lambda$.
Then
\[
 n+s=2\lambda-2m+1-g=g-1-2k_0,\qquad c=2-g+k_0.
\]
Here $s\geq1$, and $0\leq n+s\leq g-1$ follows from $n\geq0$
and $\lambda-m\leq g-1$.  Thus $0\leq k_0\leq\lfloor(g-1)/2\rfloor$.
The assertion is \eqref{pw:right-mutation}, and \eqref{pw:prefix}
identifies its mutation category with the stated existing blocks.
\end{proof}

Process the blocks with $\lambda\geq g$ using this lemma and
Corollary~\ref{pw:available-prefix}
(b).  Each is right-mutated through
the indicated initial segment and then commuted past the rest of
\eqref{dthehwe5jwjw6}.  At the end, the family
\eqref{dthehwe5jwjw6} is at the left of mega-block IV, unchanged,
and the other blocks lie to its right in $\mathcal A'$.
For example, in genus $5$ the block with $(\lambda,m)=(6,2)$ is
mutated through all three blocks
$
 \langle\theta^{-1},\theta^{-2}\cE^{\boxtimes2},
                \theta^{-3}\cE^{\boxtimes4}\rangle$.
%We do not assert that these sufficient initial segments are minimal.

\smallskip
\noindent\textit{Mega-block I.}
Its blocks are
\begin{equation}\label{sEGsrhSRhjdt}
 \left\langle Z^{\lambda+2-g}\theta^{m-\lambda-1}
           \cE^{\boxtimes(\lambda-2m)}\right\rangle_{{0\leq\lambda\leq g-2}\atop
           {0\leq m\leq\lfloor\lambda/2\rfloor}}.
\end{equation}
Retain the family of leftmost blocks with $\lambda=g-2$:
\begin{equation}\label{sEGsrhSzxfafdbRhjdt}
 \left\langle\theta^{1-g+k}\cE^{\boxtimes(g-2-2k)}
 \right\rangle_{k=\lfloor(g-2)/2\rfloor,\ldots,0}
       =L\zeta^*\mathcal B_{g-2}(1-g).
\end{equation}
%We use left mutation directly, without dualizing these tensor bundles.

\begin{lemma}\label{megablocki}
For $0\leq\lambda\leq g-3$ and
$0\leq m\leq\lfloor\lambda/2\rfloor$, let $\mathcal B$ be the
initial segment of \eqref{sEGsrhSzxfafdbRhjdt} indexed by $k\geq m$.
Then
\[
 \mathbf L_{\mathcal B}
 \left\langle Z^{\lambda+2-g}\theta^{m-\lambda-1}
              \cE^{\boxtimes(\lambda-2m)}\right\rangle
       \subset\mathcal A.
\]
\end{lemma}

\begin{proof}
Set $n=\lambda-2m$, $s=\lambda+2-g<0$, and $c=m-\lambda-1$.
Then $n-s=g-2-2m$ and $c+s=m+1-g$.
Equations~\eqref{pw:left-mutation} and \eqref{pw:prefix} give the
claim.  The bound $0\leq m\leq\lfloor(g-2)/2\rfloor$ follows
from $2m\leq\lambda\leq g-3$.
\end{proof}

Apply 
Corollary~\ref{pw:available-prefix}
(a).  Before each left mutation,
the unprocessed block commutes past the part of
\eqref{sEGsrhSzxfafdbRhjdt} indexed by $k<m$; no retained block
is changed.  Thus the nonretained part of I moves to the left into
$\mathcal A$, and \eqref{sEGsrhSzxfafdbRhjdt} remains at its right.

\smallskip
\noindent\textit{Mega-block IIa.}
This part consists of
\begin{equation}\label{sBGASFNADNM}
 \left\langle Z^{\lambda+3-g}\theta^{m-\lambda-1}
              \cE^{\boxtimes(\lambda-2m)}\right\rangle_{{g-2\leq\lambda\leq2(g-2)}\atop
        {0\leq m\leq\lfloor\lambda/2\rfloor,\ \lambda-m\leq g-2}}.
\end{equation}
Initially retain the  family of blocks
\begin{equation}\label{AWRHQERHERJQEJ}
 \left\langle Z\theta^{1-g+k}\cE^{\boxtimes(g-2-2k)}
 \right\rangle_{k=\lfloor(g-2)/2\rfloor,\ldots,0}
       =ZL\zeta^*\mathcal B_{g-2}(1-g).
\end{equation}
For a block with $\lambda\geq g-1$, put
$k_0=g-2-\lambda+m$.  Divide the block and the retained family by
$Z$ and apply \eqref{pw:right-mutation}.  Indeed,
\[
 \begin{gathered}
 n=\lambda-2m,\quad s=\lambda+2-g,\quad c=m-\lambda-1,\\
 n+s=g-2-2k_0,\qquad c=1-g+k_0.
 \end{gathered}
\]
Thus right mutation through the initial segment $k\geq k_0$ of
\eqref{AWRHQERHERJQEJ} places this block in
$Z\mathcal A'=\mathcal A$.  The cutoff is in the displayed range:
$k_0\geq0$ follows from $\lambda-m\leq g-2$, while
$g-2-2k_0=n+s\geq0$.  
Corollary~\ref{pw:available-prefix} (b), with
$t=1$, moves each processed block after the whole retained family.

We now have the two consecutive families
\eqref{sEGsrhSzxfafdbRhjdt}, \eqref{AWRHQERHERJQEJ}.
Their boundary mutation is  given by
\begin{equation}\label{pw:positive-boundary-sequence}
 0\longrightarrow\cO_M\longrightarrow Z
       \longrightarrow\cO_Z(Z)\longrightarrow0.
\end{equation}
The first map becomes an isomorphism after $R\zeta_*$:
by Lemma~\ref{pw:direct-images} both direct images are $\cO_N$,
and the map is nonzero on the open set where $\zeta$ is an
isomorphism.  Since $H^0(N,\cO_N)=\mathbb C$, it is an isomorphism
everywhere.  Consequently $R\zeta_*\cO_Z(Z)=0$.
For $F\in\mathcal B_{g-2}(1-g)$, the counit therefore identifies
\begin{equation}\label{pw:positive-boundary-mutation}
 \mathbf L_{L\zeta^*\mathcal B_{g-2}(1-g)}(ZL\zeta^*F)
       \simeq\cO_Z(Z)\otimes L\zeta^*F\in\mathcal A.
\end{equation}
This mutates the whole family \eqref{AWRHQERHERJQEJ} to the left
of \eqref{sEGsrhSzxfafdbRhjdt}.  The already processed higher IIa
blocks are also in $\mathcal A$ and can now be commuted to that
side: one orthogonality follows from adjunction and the other from
the current semiorthogonal order.  Hence all nonretained blocks of
I and IIa are to the left of \eqref{sEGsrhSzxfafdbRhjdt}, in
$\mathcal A$.

\smallskip
\noindent\textit{Mega-block IIb.}
Its kernels have the same formula as \eqref{sBGASFNADNM}, with
$0\leq\lambda\leq g-3$.  Retain the  family
\begin{equation}\label{kjsrHwrhRJegfw}
 \left\langle\theta^{2-g+k}\cE^{\boxtimes(g-3-2k)}
 \right\rangle_{k=\lfloor(g-3)/2\rfloor,\ldots,0}
       =L\zeta^*\mathcal B_{g-3}(2-g).
\end{equation}
For $\lambda\leq g-4$, apply \eqref{pw:left-mutation} with
\[
 \begin{gathered}
 n=\lambda-2m,\quad s=\lambda+3-g,\quad c=m-\lambda-1,\\
 n-s=g-3-2m,\qquad c+s=m+2-g.
 \end{gathered}
\]
The mutation category is the initial segment $k\geq m$ of
\eqref{kjsrHwrhRJegfw}.  
Corollary~\ref{pw:available-prefix} (a)
places the other IIb blocks on its left in $\mathcal A$.
They then commute past \eqref{sEGsrhSzxfafdbRhjdt}, by adjunction
and the current semiorthogonal order.  The two retained families
from I and II are unchanged and consecutive.

\smallskip
\noindent\textit{Mega-block IIIb.}
This part has kernels
\[
 Z^{\lambda+2-g}\theta^{m-\lambda}
       \cE^{\boxtimes(\lambda-2m)},\qquad
 0\leq\lambda\leq g-3,\quad0\leq m\leq\lfloor\lambda/2\rfloor.
\]
Retain the family with $\lambda=g-3$:
\begin{equation}\label{srgrhasrharhqer}
 \left\langle Z^{-1}\theta^{3-g+k}\cE^{\boxtimes(g-3-2k)}
 \right\rangle_{k=\lfloor(g-3)/2\rfloor,\ldots,0}
       =Z^{-1}L\zeta^*\mathcal B_{g-3}(3-g).
\end{equation}
For $\lambda\leq g-4$, multiply both the block and this family by
$Z$ and apply \eqref{pw:left-mutation}, with
\[
 \begin{gathered}
 n=\lambda-2m,\quad s=\lambda+3-g,\quad c=m-\lambda,\\
 n-s=g-3-2m,\qquad c+s=m+3-g.
 \end{gathered}
\]
The required initial segment in \eqref{srgrhasrharhqer} is $k\geq m$.
After tensoring back by $Z^{-1}$, the processed blocks lie to its
left in $Z^{-1}\mathcal A=\mathcal A'$.

The family \eqref{srgrhasrharhqer} is $Z^{-1}$ times the initial
segment $k\geq1$ of \eqref{dthehwe5jwjw6}.  Denote the corresponding
subcategory on $N$ by $\mathcal B_{g-3}(3-g)$ as above.  The second
boundary mutation is  induced by
\begin{equation}\label{pw:negative-boundary-sequence}
 0\longrightarrow Z^{-1}\longrightarrow\cO_M
       \longrightarrow\cO_Z\longrightarrow0.
\end{equation}
For $F\in\mathcal B_{g-3}(3-g)$, the unit of the left adjunction
gives
\begin{equation}\label{pw:negative-boundary-mutation}
 \mathbf R_{L\zeta^*\mathcal B_{g-3}(3-g)}(Z^{-1}L\zeta^*F)
       \simeq\cO_Z\otimes L\zeta^*F[-1]\in\mathcal A'.
\end{equation}
Here $R\zeta_*(Z\cO_Z)=0$ by
\eqref{pw:positive-boundary-sequence}.  Apply this right mutation
through the actual initial segment $k\geq1$, and commute the result
past the remaining $k=0$ block as in
Corollary~\ref{pw:available-prefix} (b).  After all these boundary blocks
have been processed, the lower IIIb blocks already in $\mathcal A'$
commute past the retained family \eqref{dthehwe5jwjw6}.  All
nonretained blocks from IIIb and IV lie to its right in $\mathcal A'$.

\smallskip
\noindent\textit{Mega-block IIIa.}
Its kernels are
\[
 \begin{gathered}
 Z^{\lambda+2-g}\theta^{m-\lambda}
       \cE^{\boxtimes(\lambda-2m)},\\
 g-2\leq\lambda\leq2(g-2),\quad
 0\leq m\leq\lfloor\lambda/2\rfloor,\quad\lambda-m\leq g-2.
 \end{gathered}
\]
Retain the  family
\begin{equation}\label{dthasgarharhwjw6}
 \left\langle\theta^{2-g+k}\cE^{\boxtimes(g-2-2k)}
 \right\rangle_{k=\lfloor(g-2)/2\rfloor,\ldots,0}
       =L\zeta^*\mathcal B_{g-2}(2-g).
\end{equation}
For $\lambda\geq g-1$, put $k_0=g-2-\lambda+m$ and use
\eqref{pw:right-mutation}, since
\[
 \begin{gathered}
 n=\lambda-2m,\quad s=\lambda+2-g,\quad c=m-\lambda,\\
 n+s=g-2-2k_0,\qquad c=2-g+k_0.
 \end{gathered}
\]
The initial segment $k\geq k_0$ of
\eqref{dthasgarharhwjw6} is available, by the same bounds as for IIa.
The resulting blocks lie in $\mathcal A'$ after the retained family.
They commute also past \eqref{dthehwe5jwjw6}, by adjunction and the
current semiorthogonal order.  

The decomposition now has the four retained families
\eqref{sEGsrhSzxfafdbRhjdt}, \eqref{kjsrHwrhRJegfw},
\eqref{dthasgarharhwjw6}, and \eqref{dthehwe5jwjw6}, in this order,
consecutive and unchanged.  All other blocks from I and II are on
their left in $\mathcal A$, and all other blocks from III and IV
are on their right in $\mathcal A'$.
Move the terminal group of nonretained blocks to the front by the Serre functor
$(-)\otimes\omega_M[\dim M]$.  It then lies in $\mathcal A$.
Thus we have obtained, by adjacent mutations and the specified
orthogonal commutations, a full semiorthogonal decomposition
\[
 \begin{aligned}
 D^b(M)=\bigl\langle &\mathcal K,
 L\zeta^*\mathcal B_{g-2}(1-g),
 L\zeta^*\mathcal B_{g-3}(2-g),\\
 &L\zeta^*\mathcal B_{g-2}(2-g),
 L\zeta^*\mathcal B_{g-1}(2-g)\bigr\rangle,
 \qquad \mathcal K\subset\mathcal A.
 \end{aligned}
\]
Here the second retained family is omitted when $g=2$.
For $F\in D^b(N)$, apply $R\zeta_*$ to the filtration of
$L\zeta^*F$ associated with this decomposition.  The $\mathcal K$
term vanishes, and $R\zeta_*L\zeta^*F=F$.
Hence the retained families generate $D^b(N)$.
Their full faithfulness and semiorthogonality were already available
from Theorem~\ref{LoomTheorem} and full faithfulness of $L\zeta^*$.
This proves Theorem~\ref{MainTheorem}.
\end{proof}

\end{document}
